\chapter{Progettazione e verifica dei componenti}
\label{cap:progettazione}

\section{Metodo di analisi agli elementi finiti}
\label{sec:metodo-fem}
% Esposto una volta sola; tutte le verifiche successive richiamano questa
% sezione invece di ripetersi.
%   - perche' il calcolo numerico e' necessario per queste geometrie
%   - tipi di elemento adottati e ragione della scelta
%   - discretizzazione e controllo di convergenza (da mettere per iscritto)
%   - modello di materiale e provenienza dei dati
%   - rappresentazione di vincoli, contatti e precarichi
%   - criterio di verifica e coefficiente di sicurezza assunto
%   - cio' che il modello NON rappresenta, e perche'
% Riferimento: Santi e Cesari, CLUEB 2024 (voce bibliografica da aggiungere).

\section{Scelta della tavola di posizionamento}
\label{sec:tavola}
% Selezione a catalogo, non verifica strutturale. Precede i componenti
% progettati perche' ne condiziona le interfacce.
%   - criteri dal capitolo 2: capacita' di carico, rigidezza, massa e momenti
%     parassiti, ingombro e accessibilita' delle testine, dima di foratura
%   - confronto fra le alternative considerate
%   - scelta e conseguenze sulle due piastre

\section{Clamshell di protezione dello strumento}
\label{sec:clamshell}
% Schema comune alle quattro sezioni sui componenti progettati:
% requisiti (cap. 2) -> soluzione costruttiva e materiale -> carichi e
% condizioni al contorno -> risultati della verifica -> esito.
%   - carichi: forza di serraggio delle ganasce, coppia da trasmettere
%   - verifica: pressione trasmessa al gambo, tenuta senza strisciamento
%   - esito: componente sacrificale, criterio di sostituzione

\section{Sede del blocchetto}
\label{sec:sede}
%   - carichi: coppia di reazione e peso
%   - verifica: resistenza e soprattutto rigidezza, perche' ogni cedevolezza
%     qui si somma direttamente fra provino e cella
%   - sostituzione del blocchetto fra una prova e la successiva

\section{Piastra di accoppiamento inferiore}
\label{sec:piastra-inferiore}
%   - interfacce: sei fori M2 e spina di centraggio sul lato cella,
%     dima della tavola sul lato opposto
%   - requisiti dominanti: parallelismo, eccentricita' contenuta
%   - verifica: rigidezza flessionale

\section{Piastra di accoppiamento superiore}
\label{sec:piastra-superiore}
%   - interfacce: tavola da un lato, mandrino della traversa dall'altro
%   - verifica: rigidezza

\section{Assieme e procedura di prova}
\label{sec:assieme-procedura}
%   - sequenza di montaggio
%   - centraggio manuale del blocchetto
%   - azzeramento della cella prima dell'avvio
%   - sostituzione dei componenti di consumo (clamshell, blocchetto)
